\section{Conclusion}
\label{sec:conclusion}

We have argued that a personal AI assistant can be built local-first on a consumer machine, and that doing so is a design problem in its own right rather than a scaled-down copy of the cloud. The constraints---one small GPU shared by several models, latency paid directly by one user, the assistant living next to the user's own files and credentials, and hardware that changes over time---reshape every feature. The result is an architecture organised around nine principles: local by default with governed cloud use; nothing model-based on the reply path that can run after it; stable prompt prefixes; execution over opinion; one approval model; untrusted content as data; measured and reversible change; honest, loyal, non-dependent behaviour; and hardware as an input.

The design covers the full feature set users now expect---memory, history, projects, search with citations, files, skills, a code interpreter, a canvas, connectors, images, voice---and goes beyond it to an agent company that produces tested code on one GPU, to always-on operation controlled from a phone, to self-improvement gated by local exams, and finally to calls. Each piece is specified against an 8~GB budget with its failure modes named and a staged plan with exit criteria; a build order of nine stages puts security before exposure and ends every stage with something used daily. A prototype implements the foundation and the first two stages.

What we hope readers take from this is less the feature list than the method: write down the hardware budget, put the security base before the capability that needs it, decide for every feature who holds the GPU when, and refuse to apply any change the system cannot prove on its owner's own machine. We believe this method generalises beyond \hivi{} to any assistant that is meant to be owned rather than rented.
